\clearpage
\subsection{Arkusfunktionen als Umkehrung der Winkelfunktionen}
\label{subsec:real_functions_arc_functions}
Sinus, Kosinus und Tangens wiederholen ihre Werte.
Für eindeutige Umkehrfunktionen beschränken wir sie deshalb auf je ein
Intervall, auf dem sie streng monoton sind. Die Rückwärtsfrage lautet:
Welcher Winkel gehört zu einem vorgegebenen Funktionswert?
% Handschrift: PDF-Seiten 57, 59 und 61 (jeweils gewählter Ast).
\begin{center}
  \CachedTikzFigure{content/resources/chapters/real_functions/tikz/functions_arc_restrictions.tex}
\end{center}
\par\noindent\begin{minipage}{\linewidth}
\begin{DefinitionBox}{Arkussinus, Arkuskosinus und Arkustangens}{real_functions_arc_functions}
  Die \emph{Arkusfunktionen} \EnglishTerm{inverse trigonometric functions}
  sind die Umkehrfunktionen der folgenden Einschränkungen:
  \begin{center}
    \renewcommand{\arraystretch}{1.6}
    \begin{tabular}{@{}lll@{}}
      \toprule
      Ausgangsfunktion & Gewähltes Intervall & Umkehrfunktion\\ \midrule
      $\sin$ & $[-\pi/2,\pi/2]$ &
        $\arcsin:[-1,1]\to[-\pi/2,\pi/2]$\\
      $\cos$ & $[0,\pi]$ &
        $\arccos:[-1,1]\to[0,\pi]$\\
      $\tan$ & $(-\pi/2,\pi/2)$ &
        $\arctan:\R\to(-\pi/2,\pi/2)$\\
      \bottomrule
    \end{tabular}
  \end{center}
  Sie heißen \emph{Arkussinus} \EnglishTerm{arcsine},
  \emph{Arkuskosinus} \EnglishTerm{arccosine} und
  \emph{Arkustangens} \EnglishTerm{arctangent}.
\end{DefinitionBox}
\begin{ExampleBox}[before skip=0.3\baselineskip]{Der festgelegte Winkelbereich entscheidet}{real_functions_arc_values}
  Bestimmen Sie jeweils den Winkel in der \textcolor{green!55!black}{Wertemenge} der Umkehrfunktion.
  \[
    \arcsin\frac12=\Cloze{2cm}{\frac\pi6},
    \qquad
    \arccos\left(-\frac12\right)=\Cloze{2cm}{\frac{2\pi}3},
    \qquad
    \arctan(-1)=\Cloze{2cm}{-\frac\pi4}.
  \]
\end{ExampleBox}
\end{minipage}

% Handschrift: PDF-Seiten 58, 60 und 62 (geometrische Konstruktion aller drei Umkehrungen).
Die eingeschränkten Graphen werden an $y=x$ gespiegelt.
Die blauen Ausgangsfunktionen und die grünen Umkehrfunktionen zeigen
jeweils denselben Zusammenhang mit vertauschten Koordinaten:
\begin{center}
  \CachedTikzFigure{content/resources/chapters/real_functions/tikz/functions_arc_reflections.tex}
\end{center}
\begin{ExampleBox}{Besondere Werte der Arkusfunktionen}{real_functions_arc_special_values}
  Nutzen Sie bekannte Sinus-, Kosinus- und Tangenswerte im jeweils
  gewählten Winkelbereich. Ergänzen Sie die Umkehrwerte:
  \[
    \begin{array}{c||c|c|c}
      x&-1&0&1\\\hline
      \arcsin x&\Cloze{2cm}{-\pi/2}&\Cloze{2cm}{0}&\Cloze{2cm}{\pi/2}\\[0.7em]
      \arccos x&\Cloze{2cm}{\pi}&\Cloze{2cm}{\pi/2}&\Cloze{2cm}{0}\\[0.7em]
      \arctan x&\Cloze{2cm}{-\pi/4}&\Cloze{2cm}{0}&\Cloze{2cm}{\pi/4}
    \end{array}
  \]
\end{ExampleBox}

\begin{center}
  \CachedTikzFigure{content/resources/chapters/real_functions/tikz/functions_arc_graphs.tex}
\end{center}
\begin{TheoremBox}{Eigenschaften der Arkusfunktionen}{real_functions_arc_properties}
  Alle drei Arkusfunktionen sind auf ihrer Definitionsmenge stetig.
  Arkussinus und Arkustangens sind streng monoton steigend und ungerade.
  Arkuskosinus ist streng monoton fallend.
  Alle drei sind beschränkt. Arkussinus und Arkustangens haben ihre
  einzige Nullstelle und ihren $y$-Achsenabschnitt bei $0$.
  Arkuskosinus hat seine einzige Nullstelle bei $x=1$ und den
  $y$-Achsenabschnitt $\pi/2$.
  Die Symmetrie- und Schrankenbedingungen lauten explizit:
  \[
    \arcsin(-x)=-\arcsin x\quad(x\in[-1,1]),\qquad
    \arctan(-x)=-\arctan x\quad(x\in\R),
  \]
  \[
    -\frac\pi2\le\arcsin x\le\frac\pi2,\qquad
    0\le\arccos x\le\pi\quad(x\in[-1,1]),\qquad
    -\frac\pi2<\arctan x<\frac\pi2\quad(x\in\R).
  \]

  Der Arkustangens ist beschränkt und besitzt zwei waagerechte Asymptoten:
  \[
    \lim_{x\to-\infty}\arctan x=-\frac\pi2,\qquad
    \lim_{x\to+\infty}\arctan x=\frac\pi2.
  \]
\end{TheoremBox}

\begin{KeyPointBox}{Rückwärts nur im gewählten Winkelbereich}{real_functions_arc_restrictions}
  Die beiden Richtungen der Umkehrung haben unterschiedliche Voraussetzungen:
  \[
    \begin{aligned}
      \sin(\arcsin y)&=y&&\text{für }y\in[-1,1],\\
      \arcsin(\sin x)&=x&&\text{für }x\in[-\pi/2,\pi/2].
    \end{aligned}
  \]
  Für Kosinus und Tangens lauten die Beziehungen ausdrücklich:
  \[
    \begin{aligned}
      \cos(\arccos y)&=y&&\text{für }y\in[-1,1],\\
      \arccos(\cos x)&=x&&\text{für }x\in[0,\pi],\\[0.6em]
      \tan(\arctan y)&=y&&\text{für }y\in\R,\\
      \arctan(\tan x)&=x&&\text{für }x\in(-\pi/2,\pi/2).
    \end{aligned}
  \]
\end{KeyPointBox}
\begin{ExampleBox}{Ein Umkehrwert ist noch nicht die ganze Lösungsmenge}{real_functions_trig_equation}
  Lösen Sie $\sin x=1/2$ für alle $x\in\R$.
  Nutzen Sie nach dem Arkussinus auch Symmetrie und Periodizität.
  \begin{ClozeArea}[34]{9}
    Der Arkussinus liefert zunächst $x=\pi/6$.
    Innerhalb einer Periode hat auch $\pi-\pi/6=5\pi/6$ denselben Sinus.
    Durch Wiederholung ergeben sich sämtliche Lösungen:
    \[
      x=\frac\pi6+2k\pi
      \quad\text{oder}\quad
      x=\frac{5\pi}6+2k\pi,\qquad k\in\Z.
    \]
    Deshalb ist etwa $\arcsin(\sin(5\pi/6))=\pi/6$ und nicht $5\pi/6$:
    Die Umkehrfunktion gibt immer einen Winkel in $[-\pi/2,\pi/2]$ zurück.
  \end{ClozeArea}
\end{ExampleBox}
\begin{RemarkBox}
  Die Taschenrechnertaste $\sin^{-1}$ bezeichnet meist $\arcsin$,
  nicht $1/\sin$. Prüfen Sie außerdem den Winkelmodus:
  Ein Ergebnis in Grad und eines im Bogenmaß beschreiben denselben
  Winkel mit unterschiedlichen Zahlen.
\end{RemarkBox}

% Handschrift: PDF-Seiten 63–65; Rückblick nach Einführung aller Funktionsklassen.
\subsection{Stetigkeit der bekannten Funktionen}
\label{subsec:real_functions_elementary_continuity}
Die meisten bisher behandelten Funktionen sind auf ihrer natürlichen
Definitionsmenge stetig. Die Rechenregeln aus
Satz~\ref{thm:real_functions_continuity_rules} übertragen dies auf ihre Kombinationen.
\begin{center}
  \renewcommand{\arraystretch}{1.4}
  \begin{tabular}{@{}ll@{}}
    \toprule
    Funktion & Stetig auf\\\midrule
    Polynome, $e^x$, $\sin x$, $\cos x$ & $\R$\\
    Gebrochen-rationale Funktionen $p(x)/q(x)$ & $\{x\in\R\mid q(x)\ne0\}$\\
    $\sqrt[n]{x}$, $n\ge1$ gerade & $[0,\infty)$\\
    $\sqrt[n]{x}$, $n\ge1$ ungerade & $\R$\\
    $\log_a x$, $a>0$, $a\ne1$ & $(0,\infty)$\\
    \bottomrule
  \end{tabular}
\end{center}
\begin{ExampleBox}{Stetigkeit zusammengesetzter Funktionen erkennen}{real_functions_continuous_combinations}
  Benennen Sie jeweils die Rechenregel, die ohne neue Grenzwertrechnung
  Stetigkeit auf ganz $\R$ liefert:
  \[
    \begin{array}{l@{\qquad}l}
      \sin x\,e^{-x}&\Cloze{5cm}{\text{Produkt stetiger Funktionen}}\\[0.8em]
      x^2+\cos x&\Cloze{5cm}{\text{Summe stetiger Funktionen}}\\[0.8em]
      \sin(x^2)&\Cloze{5cm}{\text{Verkettung stetiger Funktionen}}
    \end{array}
  \]
\end{ExampleBox}
Bei stückweise definierten Funktionen müssen die Übergänge gesondert
geprüft werden. Die Schaltfunktion aus
Beispiel~\ref{ex:real_functions_heaviside_continuity} zeigt, wie ein Sprung die Stetigkeit verhindert.
